% $Id: template.tex 11 2007-04-03 22:25:53Z jpeltier $

\documentclass{styling/vgtc}                          % final (conference style)
%\documentclass[review]{styling/vgtc}                 % review
%\documentclass[widereview]{styling/vgtc}             % wide-spaced review
%\documentclass[preprint]{styling/vgtc}               % preprint
%\documentclass[electronic]{styling/vgtc}             % electronic version

%% Uncomment one of the lines above depending on where your paper is
%% in the conference process. ``review'' and ``widereview'' are for review
%% submission, ``preprint'' is for pre-publication, and the final version
%% doesn't use a specific qualifier. Further, ``electronic'' includes
%% hyperreferences for more convenient online viewing.

%% Please use one of the ``review'' options in combination with the
%% assigned online id (see below) ONLY if your paper uses a double blind
%% review process. Some conferences, like IEEE Vis and InfoVis, have NOT
%% in the past.

%% Figures should be in CMYK or Grey scale format, otherwise, colour 
%% shifting may occur during the printing process.

%% These three lines bring in essential packages: ``mathptmx'' for Type 1 
%% typefaces, ``graphicx'' for inclusion of EPS figures. and ``times''
%% for proper handling of the times font family.

\usepackage{mathptmx}
\usepackage{graphicx}
\usepackage{times}
\usepackage[hidelinks]{hyperref}

%% We encourage the use of mathptmx for consistent usage of times font
%% throughout the proceedings. However, if you encounter conflicts
%% with other math-related packages, you may want to disable it.

%% If you are submitting a paper to a conference for review with a double
%% blind reviewing process, please replace the value ``0'' below with your
%% OnlineID. Otherwise, you may safely leave it at ``0''.
\onlineid{0}

%% declare the category of your paper, only shown in review mode
\vgtccategory{Research}

%% allow for this line if you want the electronic option to work properly
\vgtcinsertpkg

%% In preprint mode you may define your own headline.
%\preprinttext{To appear in an IEEE VGTC sponsored conference.}

%% This is how authors are specified in the conference style

%% Author and Affiliation (single author).
%%\author{Roy G. Biv\thanks{e-mail: roy.g.biv@aol.com}}
%%\affiliation{\scriptsize Allied Widgets Research}

%% Author and Affiliation (multiple authors with single affiliations).
%%\author{Roy G. Biv\thanks{e-mail: roy.g.biv@aol.com} %
%%\and Ed Grimley\thanks{e-mail:ed.grimley@aol.com} %
%%\and Martha Stewart\thanks{e-mail:martha.stewart@marthastewart.com}}
%%\affiliation{\scriptsize Martha Stewart Enterprises \\ Microsoft Research}

%% Author and Affiliation (multiple authors with multiple affiliations)
\title{DreamGuard - Checkpoint 2}

\author{Andrew Scouten\\ %
        \scriptsize{yzb2@txstate.edu}\\ %
        \scriptsize Texas State University
\and Avery R Vanausdal\\ %
        \scriptsize{arv107@txstate.edu}\\ %
        \scriptsize Texas State University}

%% A teaser figure can be included as follows, but is not recommended since
%% the space is now taken up by a full width abstract.
%\teaser{
%  \includegraphics[width=1.5in]{styling/sample.eps}
%  \caption{Lookit! Lookit!}
%}
\begin{document}

\maketitle
\thispagestyle{empty}
\vspace{-1em}

\section*{Current Status}

DreamGuard has advanced from a single-condition proof-of-concept to a
fully-instrumented multi-condition platform ready for a within-subject
evaluation. The two unresolved items from Checkpoint~1---the per-eye
double-vision artifact and the absence of additional blending styles---are
both resolved. All study conditions are implemented, proximity detection is
automated, and session logging is in place.

\section*{What Has Been Implemented}

\textbf{Multi-condition transition system.}
A new \texttt{DreamGuardTransition} node provides four spatially-rendered
fullscreen overlay styles via a two-pass shader architecture
(a \texttt{blend\_add} ceiling-fix pass followed by a \texttt{blend\_mul}
pattern pass). The five study conditions now available are:
\textit{NONE} (control, no warning),
\textit{GRID} (3-D perspective guardian grid --- Meta baseline),
\textit{FOG\_BLEND} (peripheral fog that reveals passthrough, condition~A),
\textit{BLOOM\_BLEND} (warm bloom from edges with glowing border, condition~B), and
\textit{PERIPHERAL\_VIGNETTE} (passthrough rim with clear VR centre, condition~C).
The fragment-dissolve and rectangular-window modes from Checkpoint~1 are
retained as additional conditions.

\textbf{World-space guardian grid shader.}
The per-eye stereo rendering artifact identified at Checkpoint~1 is
eliminated. The grid shader reconstructs world-space geometry via ray--plane
intersection (floor and ceiling planes parameterised by \texttt{floor\_y} and
\texttt{room\_height}), avoiding the \texttt{SCREEN\_UV} stereo-target
ambiguity entirely. Grid lines are blue-tinted; gaps are transparent (ALPHA=0),
exposing passthrough. A distance fade dissolves the grid beyond 8--22~m so
the horizon remains uncluttered.

\textbf{Orchestrator and API.}
A \texttt{DreamGuard} orchestrator node centralises style switching, XR blend
mode management (\texttt{transparent\_bg}, \texttt{XR\_ENV\_BLEND\_MODE\_ALPHA\_BLEND}),
and the trigger/clear/velocity-scaled-trigger API. Styles can be cycled
programmatically (\texttt{next\_style()}) or set directly, allowing
counterbalanced condition ordering during the study.

\textbf{Automated proximity detection.}
\texttt{DreamGuardProximity} casts $N$ evenly-spaced horizontal rays from
approximately waist height each physics frame. Trigger intensity scales
proportionally with proximity ($0$ at the warning threshold, $1$ at contact),
so the overlay strengthens as the player approaches the boundary.
A velocity-scaled variant (motivated by Google's GAP findings on locomotion
speed and cybersickness) is also available.

\textbf{Session instrumentation.}
\texttt{DreamGuardLogger} connects to the orchestrator's signals and writes a
timestamped CSV to on-device storage with six event types:
\textsc{session\_start}, \textsc{session\_end}, \textsc{style\_change},
\textsc{trigger}, \textsc{clear}, and \textsc{blend\_sample} (configurable
interval, default 0.5~s). Logs are retrievable via \texttt{adb pull} after
each session.

\section*{What Remains}

\begin{enumerate}
  \item \textbf{IRB and study protocol.} The within-session SSQ carryover
    concern noted at Checkpoint~1 must be addressed (e.g.\ washout periods or
    a between-subjects design for the SSQ measure). The no-warning control
    condition should be reframed as a minimal-cue baseline to reduce IRB risk.
    Target IRB submission and approval before data collection begins.

  \item \textbf{Data collection.} Recruit participants for the within-subject
    evaluation. Each session will counterbalance condition order and collect
    SSQ (pre/post per condition), IPQ presence ratings, and the objective
    logged near-miss counts. Pilot testing with $\ge 3$ participants before
    the full run to validate the protocol and logging pipeline.

  \item \textbf{Analysis.} Repeated-measures ANOVA or Friedman test (depending
    on SSQ distribution) across conditions for each dependent variable.
    Pairwise comparisons of DreamGuard conditions against the GRID baseline
    and NONE control. Logged blend-sample time series will be summarised per
    condition to characterise intervention duration and intensity.

  \item \textbf{Final write-up.} Consolidate results, limitations (ecological
    validity of the dungeon environment, sample size), and directions for
    future work (dynamic style selection, adaptive threshold calibration).
\end{enumerate}

\end{document}
